\subsection{有限线性群的不变量：环论性质}

定理 3. 我们沿用定理 1 的假设与记号。在 \(R_{+}\) （R 的理想）的由齐次元素组成的生成元组的集合中，取一个极小元 \((\alpha_{1},\ldots,\alpha_{l})\) 。设 \(k_{i}\) 是 \(\alpha_{i}\) 的次数。假设诸 \(k_{i}\) 与 K 的特征指数互素。则 \(l=\dim V\) ，诸 \(\alpha_{i}\) 生成 K-代数 R，且在 K 上代数无关。特别地，R 是 K 上超越次数为 l 的一个分次多项式 K-代数。

关于 \(k_{i}\) 所作的假设是多余的，但对有限 Coxeter 群的应用而言无关紧要，因为那时 K = R。参见第 5 号，我们将在那里给出定理 3 的另一个证明。

定理 3 由命题 2 (i)、定理 1 以及下面的引理得出：

引理 1. 设 K 是一个交换域，S 是一个分次多项式 K-代数，R 是 S 的一个有限型分次子代数，使得 R-模 S 具有一个由齐次元素组成的基 \((z_{\lambda})_{\lambda \in \Lambda}\) 。在 \(R_{+}\) （R 的理想）的由齐次元素组成的生成元组的集合中，取一个极小元 \((\alpha_{1}, \ldots, \alpha_{s})\) 。假设对所有 i，次数 \(k_{i}\) （\(\alpha_{i}\) 的）与 K 的特征指数 p 互素。则诸 \(\alpha_{i}\) 生成 K-代数 R 且在 K 上代数无关。

由《代数》第 II 章 § 11 第 4 号命题 7，关于 \(\alpha_{i}\) 所作的假设等价于说它们是齐次的，且它们的像在 K-向量空间 \(R_{+}/(R_{+})^{2}\) 中构成该空间的基。这一条件在基域扩张下不变；于是我们可以化归到基域为完满的情形。

族 \((\alpha_{1},\ldots,\alpha_{s})\) 生成代数 R（《交换代数》第 III 章，§1，第 2 号，命题 1）。我们用反证法，假设这个族在 K 上不是代数无关的。

\begin{enumerate}
\def\labelenumi{\arabic{enumi})}
\tightlist
\item
  我们首先证明存在诸族
\end{enumerate}

\[
\left(\beta_ {i}\right) _ {1 \leqslant i \leqslant s}, \quad \left(y _ {k}\right) _ {1 \leqslant k \leqslant r}, \quad \left(d _ {i k}\right) _ {1 \leqslant i \leqslant s, 1 \leqslant k \leqslant r}
\]

它们由 S 的齐次元素组成，且满足下列性质：

\[
\beta_ {i} \in R \text {for all} i, \text {and the} \beta_ {i} \text {are not all zero};\tag{7}
\]

\[
\deg y _ {k} > 0 \text {   for   all   } k;\tag{8}
\]

\[
\alpha_ {i} = \sum_ {k = 1} ^ {r} d _ {i k} y _ {k} \quad \text { for   all } i;\tag{9}
\]

\[
\sum_ {i = 1} ^ {s} \beta_ {i} d _ {i k} = 0 \quad \text { for   all } k.\tag{10}
\]

设 \(X_{1},\ldots,X_{s}\) 是未定元，并给 \(K[X_{1},\ldots,X_{s}]\) 赋予使 \(X_{i}\) 的次数为 \(k_{i}\) 的分次代数结构。存在非零齐次元素 \(H(X_{1},\ldots,X_{s})\) ，它属于 \(K[X_{1},\ldots,X_{s}]\) ，使得 \(H(\alpha_{1},\ldots,\alpha_{s})=0\) ；取 H 为次数最小的。若 \(\partial H/\partial X_{i}\neq0\) ，则多项式 \(\frac{\partial H}{\partial X_{i}}(\alpha_{1},\ldots,\alpha_{s})\) 是 R 的一个非零齐次元素；若 \(p\neq1\) ，则 H 不具有 \(H_{1}^{p}\) （其中 \(H_{1}\in K[X_{1},\ldots,X_{s}]\) ）的形式。取

\[
\beta_ {i} = k _ {i} \frac {\partial \mathrm{H}}{\partial \mathrm{X} _ {i}} (\alpha_ {1}, \dots , \alpha_ {s}).
\]

由于 K 是完满的，多项式 \(\partial H/\partial X_{i} \in K[X_{1},\ldots,X_{s}]\) 不全为零（《代数》第 V 章，§1，第 3 号，命题 4）；鉴于关于 \(k_{i}\) 的假设，诸 \(\beta_{i}\) 亦不全为零。

另一方面，S 可与一个分次多项式代数 \(K[x_{1},\ldots,x_{r}]\) 等同，其中 \(x_{1},\ldots,x_{r}\) 是适当的未定元，次数为适当的 \(m_{i}>0\) 。设 \(D_{k}\) 是 S 上关于 \(x_{k}\) 的偏导数。取 \(d_{ik} = k_i^{-1}\mathrm{D}_k(\alpha_i)\) 。则等式 (10) 成立，因为其左边是 \(\mathrm{D}_k(\mathrm{H}(\alpha_1,\dots ,\alpha_s))\) 。另一方面，若置 \(y_{1} = m_{1}x_{1},\ldots ,y_{r} = m_{r}x_{r}\) ，则等式 (9) 由第 1 号的等式 (5) 得出。

\begin{enumerate}
\def\labelenumi{\arabic{enumi})}
\setcounter{enumi}{1}
\tightlist
\item
  设 \(\mathfrak{b}\) 是 \(R\) 的由诸 \(\beta_{i}\) 生成的理想；存在下述集合的一个子集 \(J\) ：
\end{enumerate}

\[
\mathrm{I} = \{1, 2, \dots , s \}
\]

使得 \((\beta_{i})_{i\in\mathbb{J}}\) 是 b 的一个极小生成元组。则 \(J\neq\varnothing\) ，因为 \(b\neq0\) 。我们将由 (9) 与 (10) 推出：若 \(i\in J\) ，则 \(\alpha_{i}\) 是诸 \(\alpha_{j}\) （\(j\neq i\) ）的 R-线性组合，这与 \((\alpha_{1},\ldots,\alpha_{s})\) 的极小性矛盾，从而完成证明。

存在 R 的齐次元素 \(\gamma_{ji}\) ( \(i \in \mathrm{J}, j \in \mathrm{I} - \mathrm{J}\) ) 使得

\[
\beta_ {j} = \sum_ {i \in J} \gamma_ {j i} \beta_ {i} (j \in I - J).\tag{11}
\]

考虑到 (11)，公式 (10) 给出

\[
\sum_ {i \in \mathrm{J}} \beta_ {i} \left(d _ {i k} + \sum_ {j \in \mathrm{I} - \mathrm{J}} \gamma_ {j i} d _ {j k}\right) = 0.\tag{12}
\]

置

\[
u _ {i k} = d _ {i k} + \sum_ {j \in I - J} \gamma_ {j i} d _ {j k}.\tag{13}
\]

于是

\[
\sum_ {i \in J} \beta_ {i} u _ {i k} = 0.\tag{14}
\]

写 \(u_{ik} = \sum_{\lambda \in \Lambda} \delta_{ik\lambda} z_\lambda\) ，其中诸 \(\delta_{ik\lambda}\) 属于 R。关系式 (14) 蕴含 \(\sum_{i \in J} \beta_i \delta_{ik\lambda} = 0\) （对所有 \(k\) 与 \(\lambda\) ）。若某个 \(\delta_{ik\lambda}\) 有一个次数为 0 的非零齐次分量，则前面的等式将蕴含某个 \(\beta_i\) ( \(i \in J\) ) 是其余元素的 R-线性组合，与 \((\beta_i)_{i \in J}\) 的极小性矛盾。于是 \(\delta_{ik\lambda} \in R_+\) ，从而 \(u_{ik} \in R_+ S\) （对所有 \(i\) 与 \(k\) ）。于是存在 \(u_{ikh} \in S\) 使得 \(u_{ik} = \sum_{h=1}^{s} u_{ikh} \alpha_h\) ，换句话说，由 (13)，

\[
d _ {i k} + \sum_ {j \in I - J} \gamma_ {j i} d _ {j k} = \sum_ {h = 1} ^ {s} u _ {i k h} \alpha_ {h}.\tag{$(15)_{ik}$}
\]

把 \((15)_{ik}\) 的两边乘以 \(y_{k}\) ，并对固定的 i ∈ J 与 \(k = 1, 2, \ldots, r\) 相加；鉴于 (9)，我们得到

\[
\alpha_ {i} + \sum_ {j \in I - J} \gamma_ {j i} \alpha_ {j} = \sum_ {h = 1} ^ {s} \sum_ {k = 1} ^ {r} u _ {i k h} y _ {k} \alpha_ {h}.
\]

取两边的次数为 \(k_i\) 的齐次分量。由于 \(\deg y_k > 0\) ，\(\alpha_i\) 是诸 \(\alpha_j\) （\(j \neq i\) ）的 S-线性组合。由于 S 在 R 上自由且 \(\alpha_1, \ldots, \alpha_s \in \mathrm{R}\) ，由此可知 \(\alpha_i\) 是诸 \(\alpha_j\) （\(j \neq i\) ）的 R-线性组合（《交换代数》第 I 章，§3，第 5 号，命题 9 d))。

推论. 在定理 3 的假设与记号下，R 的特征次数之积是 Card(G)。

事实上，\(\operatorname{rk}_{\mathrm{R}}(S) = \operatorname{Card}(G)\) （公式 (6)，第 2 号）。S 的特征次数都等于 1。于是推论由第 1 号命题 2 (iii) 得出。

引理 2. 设 K 是一个交换域，V 是 K 上一个有限维向量空间，\(S = \bigoplus_{n \geqslant 0} S_n\) 是 V 的对称代数，s 是 V 的一个自同态，\(s^{(n)}\) 是 s 到 \(S_n\) 的典范开拓。则当 T 表示一个未定元时，我们在 K{[}{[}T{]}{]} 中有

\[
\sum_ {n = 0} ^ {\infty} \operatorname{Tr} (s ^ {(n)}) \mathrm{T} ^ {n} = (\det (1 - s \mathrm{T})) ^ {- 1}.
\]

必要时扩张基域，我们可以假设 K 是代数闭的。设 \((e_{1},\ldots,e_{r})\) 是 V 的一个基，s 关于它的矩阵是下三角的，并设 \(\lambda_{1},\ldots,\lambda_{r}\) 是这个矩阵的对角元素。关于按字典序排列的基 \((e_{1}^{i(1)}\ldots e_{r}^{i(r)})_{i(1)+\cdots+i(r)=n}\) （\(S_{n}\) 的），\(s^{(n)}\) 的矩阵是下三角的，其对角元素是诸乘积 \(\lambda_{1}^{i(1)}\ldots\lambda_{r}^{i(r)}\) 。于是

\[
\operatorname{Tr} \left(s ^ {(n)}\right) = \sum_ {i (1) + \dots + i (r) = n} \lambda_ {1} ^ {i (1)} \dots \lambda_ {r} ^ {i (r)},
\]

从而

\[
\begin{array}{l} \sum_ {n = 0} ^ {\infty} \operatorname{Tr} (s ^ {(n)}) \mathrm{T} ^ {n} = \left(\sum_ {n = 0} ^ {\infty} \lambda_ {1} ^ {n} \mathrm{T} ^ {n}\right) \left(\sum_ {n = 0} ^ {\infty} \lambda_ {2} ^ {n} \mathrm{T} ^ {n}\right) \ldots \left(\sum_ {n = 0} ^ {\infty} \lambda_ {r} ^ {n} \mathrm{T} ^ {n}\right) \\ \qquad = (1 - \lambda_ {1} \mathrm{T}) ^ {- 1} (1 - \lambda_ {2} \mathrm{T}) ^ {- 1} \ldots (1 - \lambda_ {r} \mathrm{T}) ^ {- 1} \\ \qquad = (\det (1 - s \mathrm{T})) ^ {- 1}. \end{array}
\]

引理 3. 设 K、V 与 S 如引理 2 中所述，G 是 V 的一个有限自同构群，q 是 G 的阶，R 是 S 中由在 G 下不变的元素组成的分次子代数。假设 K 的特征为 0。则 R 的 Poincaré 级数是

\[
q ^ {- 1} \sum_ {g \in G} (\det (1 - g T)) ^ {- 1}.
\]

事实上，自同态 \(f = q^{-1} \sum_{g \in G} g^{(n)}\) 是 \(S_n\) 到 \(R_n\) 上的一个投影，故 \(\operatorname{Tr}(f) = \dim_K S_n^G\) 。于是 \(R\) 的 Poincaré 级数是

\[
q ^ {- 1} \sum_ {g \in G} \left(\sum_ {n = 0} ^ {\infty} (\operatorname{Tr} g ^ {(n)}) T ^ {n}\right),
\]

于是只需应用引理 2。

命题 3. 在定理 3 的假设与记号下，设 H 是属于 G 且异于 1 的伪反射所成的集合。假设 K 的特征为 0。则 \(\operatorname{Card}(\mathrm{H}) = \sum_{i=1}^{l}(k_i - 1)\) 。

由附录命题 3，我们可以假设 K 是代数闭的。对任意 \(g \in G\) ，设 \(\lambda_{1}(g), \ldots, \lambda_{l}(g)\) 是它的特征值。由于每个 \(g \in G\) 都可对角化（附录，命题 2），g = 1 当且仅当所有 \(\lambda_{i}(g)\) 都等于 1，而 \(g \in H\) 当且仅当等于 1 的 \(\lambda_{i}(g)\) 的个数是 l - 1（此时我们用 \(\lambda(g)\) 表示异于 1 的特征值）。第 1 号的命题 1 与引理 3 证明

\[
q \prod_ {i = 1} ^ {l} (1 - T ^ {k _ {i}}) ^ {- 1} = \sum_ {g \in G} (\det (1 - g T)) ^ {- 1}\tag{16}
\]

在 \(\mathbf{K}[[\mathbf{T}]]\) 中成立，从而在 \(\mathbf{K}(\mathbf{T})\) 中成立。于是，我们在 \(\mathbf{K}(\mathbf{T})\) 中有

\[
q \frac {(1 - T) ^ {l - 1}}{\prod_ {i = 1} ^ {l} (1 - T ^ {k _ {i}})} = \frac {1}{1 - T} + \sum_ {g \in H} \frac {1}{1 - \lambda (g) T} + \sum_ {g \neq 1, g \notin H} \frac {(1 - T) ^ {l - 1}}{\det (1 - g T)},
\]

它可以写成

\[
\begin{array}{c} \frac {q - \prod_ {i = 1} ^ {l} (1 + T + \cdots + T ^ {k _ {i} - 1})}{(1 - T) \prod_ {i = 1} ^ {l} (1 + T + \cdots + T ^ {k _ {i} - 1})} \\ = \sum_ {g \in H} \frac {1}{1 - \lambda (g) T} + \sum_ {g \neq 1, g \notin H} \frac {(1 - T) ^ {l - 1}}{\det (1 - g T)}. \end{array}\tag{17}
\]

由此可知 \(q - \prod_{i=1}^{l} (1 + T + \cdots + T^{k_i - 1})\) 在 T = 1 处为零，故 \(q = k_1 k_2 \ldots k_l\) ，这一点我们已由定理 3 的推论得知。鉴此，设 \(Q(T)\) 是多项式 \((1 - T)^{-1}(q - \prod_{i=1}^{l} (1 + T + \cdots + T^{k_i - 1}))\) 。对等式 \((1 - T)Q(T) = q - \prod_{i=1}^{l} (1 + T + \cdots + T^{k_i - 1})\) 求导并置 T = 1，我们看到 \(-Q(1)\) 是下式在 T = 1 处的值

\[
\begin{array}{l} - \frac {d}{d t} \left(\prod_ {i = 1} ^ {l} (1 + T + \dots + T ^ {k _ {i} - 1})\right) \\ = - \sum_ {i = 1} ^ {l} \left((1 + 2 T + \dots + (k _ {i} - 1) T ^ {k _ {i} - 2}) \prod_ {j \neq i} (1 + T + \dots + T ^ {k _ {j} - 1})\right) \end{array}
\]

于是

\[
Q (1) = \sum_ {i = 1} ^ {l} \frac {\left(k _ {i} - 1\right) k _ {i}}{2} \prod_ {j \neq i} k _ {j} = \left(\prod_ {j = 1} ^ {l} k _ {j}\right) \left(\sum_ {i = 1} ^ {l} \frac {k _ {i} - 1}{2}\right).
\]

回到 (17)，我们另一方面有

\[
Q (1) = \left(\prod_ {j = 1} ^ {l} k _ {j}\right) \left(\sum_ {g \in H} \frac {1}{1 - \lambda (g)}\right).
\]

于是，

\[
\sum_ {i = 1} ^ {l} \frac {k _ {i} - 1}{2} = \sum_ {g \in H} \frac {1}{1 - \lambda (g)}.\tag{18}
\]

现在，\(\mathbf{G}\) 中保持一个给定超平面的点固定的元素保持该超平面的补直线稳定（附录，命题 2），从而由《代数》第 V 章 §11 第 1 号定理 1 构成一个循环子群 \(\mathbf{G}'\) （\(\mathbf{G}\) 的）。设 \(t\) 是 \(\mathbf{G}'\) 的阶；\(\lambda(g)\) （\(g \in \mathbf{G}'\) ）的值是 \(\theta, \theta^2, \ldots, \theta^{t-1}\) ，其中 \(\theta\) 是一个本原 \(t\) 次单位根。我们有 \(\frac{1}{1-\theta^i} + \frac{1}{1-\theta^{t-i}} = 1\) ，故 \(\sum_{g \in \mathbf{G}', g \neq 1} \frac{1}{1-\lambda(g)} = \frac{1}{2}(t-1) = \frac{1}{2}\mathrm{Card}(\mathrm{H} \cap \mathrm{G}')\) 。于是等式 (18) 证明了命题。

注. 当 \(\mathbf{K} = \mathbf{R}\) 、\(\mathbf{G}\) 是一个 Coxeter 群且 \(\mathbf{H}\) 是属于 \(\mathbf{G}\) 的反射所成的集合时，我们知道（§ 3）\(\mathbf{H}\) 的元素与 \(\mathbf{V}\) 的墙一一对应。

命题 4. 在定理 3 的假设与记号下，假设 K 的特征 \(\neq 2\) 。则 \(-1 \in G\) 当且仅当 R 的特征次数 \(k_{1}, \ldots, k_{l}\) 都是偶数。

设 \(f\) 是代数 \(S\) 的自同构，它开拓 \(-1\) （\(V\) 的自同构）。则 \(f(z) = (-1)^{\deg z}z\) （对所有齐次 \(z\) ，\(S\) 中的）。于是，若 \(-1 \in G\) ，则 \(R\) 的每个奇次数的齐次元素为零，且诸 \(k_i\) 都是偶数。反之，若诸 \(k_i\) 都是偶数，则 \(R\) 的每个元素在 \(f\) 下不变，且 Galois 理论表明 \(-1 \in G\) 。

